\documentclass[11pt]{article}

\usepackage[margin=25mm]{geometry}

\usepackage{amsmath,amssymb,fontspec,unicode-math}

\IfFontExistsTF{DejaVu Serif}{\setmainfont{DejaVu Serif}}{\setmainfont{Latin Modern Roman}}

\IfFontExistsTF{Latin Modern Math}{\setmathfont{Latin Modern Math}}{\setmathfont{DejaVu Math TeX Gyre}}

\usepackage{xurl,hyperref}

\setlength{\parindent}{0pt}\setlength{\parskip}{6pt}

\begin{document}

\title{Attention (as Discrete-Time Markov) Chains}

\author{ELVIA preliminary AI analysis}\date{21 October 2026}\maketitle

Preliminary assessment: Minor.\par

Independent scientific review remains pending. Author adoption has not been requested.

Identity addition need not slow mixing, and eigenvalue modulus alone does not certify metastability. A nonuniform personalization helper is nonstochastic; the central stationary-chain construction and empirical results remain separate.

\section*{Review scope}

Full supplied technical report read alongside the corresponding original-paper claim passages and applicable author-release source. This pass independently checks the stated reasoning and bounds; historical executions were not rerun.

\section*{Problem 1: Adding identity transitions can accelerate mixing}

Source: Section 4.1, identity-addition interpretation

Assessment: supported; impact: minor

\subsection*{Source claim (paraphrase)}

The paper states that \((I+A)/2\) preserves stationary state and slows convergence because of increased self-transition.

\subsection*{Analysis}

Take positive stochastic attention \(A=\left(\begin{smallmatrix}.1&.9\\.9&.1\end{smallmatrix}\right)\), realizable by finite logits \(Q=I,K=\sqrt2(\log A)^T\). Its nontrivial eigenvalue is \(-.8\); the lazy chain has \(.1\). From a single-state start, total variation to \((1/2,1/2)\) is \((.8)^k/2\) versus \((.1)^k/2\). Thus identity addition accelerates this already aperiodic chain. Stationary-distribution preservation remains valid.

\subsection*{Finding and scope}

The unrestricted mixing-direction explanation is false. TokenRank and reported empirical improvements are not thereby refuted.

\subsection*{Proposed correction}

State the eigenvalue map \(\lambda\mapsto(1+\lambda)/2\); mixing direction depends on the spectrum. A reversible nonnegative-spectrum restriction supports the narrower spectral-slowdown statement.

The described correction remains a proposal; author adoption or a corresponding erratum is not established.

\section*{Problem 2: Eigenvalue modulus does not certify metastable regions}

Source: Sections 3.1/4.4; Supplement E.1; dino.py eigenvalue weighting

Assessment: supported; impact: minor

\subsection*{Source claim (paraphrase)}

The magnitude of the second eigenvalue is associated with stable metastable regions and used for head weighting.

\subsection*{Analysis}

Positive stochastic matrices with diagonal entries \(.99\) or \(.01\), respectively, have nontrivial eigenvalues \(.98\) and \(-.98\), hence identical modulus. In a two-state chain, expected singleton residence time is \(1/(1-A_{ii})\): \(100\) for the first, only \(1/.99\) for the second. Large negative modulus describes persistent alternation rather than prolonged residence; the scalar also loses multiplicity information.

\subsection*{Finding and scope}

Modulus alone neither certifies nor counts metastable regions. The measured weighting correlation can still hold for observed attention heads; their negative-mode frequency is unmeasured.

\subsection*{Proposed correction}

Present modulus as an empirical spectral feature; stronger semantic-metastability claims require structural assumptions and retention/exit measurements.

The described correction remains a proposal; author adoption or a corresponding erratum is not established.

\section*{Problem 3: Nonuniform personalization breaks row stochasticity}

Source: helpers.py, prep\_for\_power\_method, lines 63–83

Assessment: supported; impact: minor

\subsection*{Source claim (paraphrase)}

The helper promises a stochastic personalized transition but uses \(E=ee^T/n\).

\subsection*{Analysis}

For \(A=\left(\begin{smallmatrix}.8&.2\\.3&.7\end{smallmatrix}\right)\), \(\alpha=.8\) and probability vector \(e=(.75,.25)^T\), the returned row sums are \(.875\) and \(.825\). Generally they equal \(\alpha+(1-\alpha)e_i\sum_je_j/n\), not one. Normalizing later power iterates produces a different Perron iteration, not a stochastic matrix. The default all-ones branch is valid.

\subsection*{Finding and scope}

The optional personalized release branch violates its contract. No evidence shows published demos used that branch.

\subsection*{Proposed correction}

For normalized positive \(p\), use \(T=\alpha A+(1-\alpha)\mathbf1p^T\). Rows sum to one; positivity gives ergodicity. Zero entries in \(p\) need additional assumptions for that stronger guarantee.

The described correction remains a proposal; author adoption or a corresponding erratum is not established.

\section*{Limits}

Preliminary AI source/argument assessment; independent human scientific adjudication remains pending.\par

No author-provided or retained verification program was executed. Published empirical results were not independently reproduced.\par

\section*{Source identity}

NeurIPS 2025 official proceedings PDF (exact retained local copy inspected)

\url{https://proceedings.neurips.cc/paper_files/paper/2025/file/4e8f257e054abd24c550d55e57cec274-Paper-Conference.pdf}

Reviewed PDF SHA-256: \nolinkurl{3eb0c58c71f2ce9ac0f7bdb00cb20290df43425a833d6b8cefacfcdc0b905903}

The SHA-256 identifies the exact retained PDF inspected in this rapid source/argument check. Original supplied report bytes are retained locally. No unambiguous original PDF retrieval timestamp was recovered; file modification time was not substituted.

\end{document}
